% GENERATED FILE. Edit canonical lessons, then run npm run generate:books.
% canonical-source-hash: fnv1a64:6abd3162250d6a4b
% canonical-lessons: BN-C155-kobji, BN-C155-phusphus, BN-C155-kidni, BN-C155-byam, BN-C155-svasthyo

\chapter{Wrist, Lung, Kidney, Exercise, Health}
\label{ch:bn-wrist-lung-kidney-exercise-health}
\hlchaptermodality{155}

\begin{chapteropening}
\textbf{By the end of this chapter:} \emph{I can say a wrist, a lung, a kidney, exercise and health.}

Complete the last lesson of chapter 155: I can say a wrist, a lung, a kidney, exercise and health.
\end{chapteropening}

\section[\texorpdfstring{\bn{কবজি} (kôbji)}{কবজি (kôbji)}]{\bn{কবজি} (kôbji) — a wrist}

\label{lesson:BN-C155-kobji}



\begin{quote}
Before the new one: say the Bengali for a headache, then the Bengali for a stomach ache.
\end{quote}

\subsection*{You'll want to know: \bn{কবজি}}
\textbf{\bn{কবজি}} — \emph{kôbji} — \textquotedblleft{}a wrist\textquotedblright{}.

\begin{cousinweb}
The origin of \textbf{\bn{কবজি}} (\emph{kôbji}) is not settled: CDIAL has no entry for it, and a search by form and by the gloss \textquotedblleft{}wrist\textquotedblright{} finds only Bengali \emph{rākhi}, \textquotedblleft{}thread worn round the wrist\textquotedblright{} (CDIAL 10551), a different word.
\end{cousinweb}

\begin{grammarlens}[title={What you've built}]
One more word for this chapter.
\end{grammarlens}

\subsection*{Your turn}
\begin{itemize}
\raggedright
  \item \emph{Say it:} \emph{kôbji}
  \item \emph{Say it:} \emph{kôbji}, once more
  \item \emph{Say it:} say \emph{kôbji}
  \item \emph{Recall:} say the Bengali for a headache, then the Bengali for a stomach ache, then say \emph{kôbji} again
\end{itemize}

\begin{tcolorbox}[breakable,colback=teal!4,colframe=teal!35!black,title={Before you move on}]
What does \bn{কবজি} mean? (\textquotedblleft{}A wrist\textquotedblright{}.) Say it once more.
\end{tcolorbox}
\section[\texorpdfstring{\bn{ফুসফুস} (phusphus)}{ফুসফুস (phusphus)}]{\bn{ফুসফুস} (phusphus) — a lung}

\label{lesson:BN-C155-phusphus}



\begin{quote}
Before the new one: say the Bengali for a stomach ache, then the Bengali for a wrist.
\end{quote}

\subsection*{You'll want to know: \bn{ফুসফুস}}
\textbf{\bn{ফুসফুস}} — \emph{phusphus} — \textquotedblleft{}a lung\textquotedblright{}.

\begin{cousinweb}
\textbf{\bn{ফুসফুস}} (\emph{phusphus}) matches Sanskrit \emph{phupphusa}, \textquotedblleft{}lungs\textquotedblright{}, which CDIAL calls probably imitative; the Bengali word inherited from it is \emph{phõprā}, \textquotedblleft{}lungs\textquotedblright{} (CDIAL 9090). CDIAL also lists a Bengali \emph{phusphus}, \textquotedblleft{}whispering\textquotedblright{}, under \emph{*phuss}, \textquotedblleft{}hissing\textquotedblright{} (CDIAL 9099).
\end{cousinweb}

\begin{grammarlens}[title={What you've built}]
One more word for this chapter.
\end{grammarlens}

\subsection*{Your turn}
\begin{itemize}
\raggedright
  \item \emph{Say it:} \emph{phusphus}
  \item \emph{Say it:} \emph{phusphus}, once more
  \item \emph{Say it:} say \emph{phusphus}
  \item \emph{Recall:} say the Bengali for a stomach ache, then the Bengali for a wrist, then say \emph{phusphus} again
\end{itemize}

\begin{tcolorbox}[breakable,colback=teal!4,colframe=teal!35!black,title={Before you move on}]
What does \bn{ফুসফুস} mean? (\textquotedblleft{}A lung\textquotedblright{}.) Say it once more.
\end{tcolorbox}
\section[\texorpdfstring{\bn{কিডনি} (kiḍni)}{কিডনি (kiḍni)}]{\bn{কিডনি} (kiḍni) — a kidney}

\label{lesson:BN-C155-kidni}



\begin{quote}
Before the new one: say the Bengali for a wrist, then the Bengali for a lung.
\end{quote}

\subsection*{You'll want to know: \bn{কিডনি}}
\textbf{\bn{কিডনি}} — \emph{kiḍni} — \textquotedblleft{}a kidney\textquotedblright{}.

\begin{cousinweb}
\textbf{\bn{কিডনি}} (\emph{kiḍni}) is a loan from English \emph{kidney}.
\end{cousinweb}

\begin{grammarlens}[title={What you've built}]
One more word for this chapter.
\end{grammarlens}

\subsection*{Your turn}
\begin{itemize}
\raggedright
  \item \emph{Say it:} \emph{kiḍni}
  \item \emph{Say it:} \emph{kiḍni}, once more
  \item \emph{Say it:} say \emph{kiḍni}
  \item \emph{Recall:} say the Bengali for a wrist, then the Bengali for a lung, then say \emph{kiḍni} again
\end{itemize}

\begin{tcolorbox}[breakable,colback=teal!4,colframe=teal!35!black,title={Before you move on}]
What does \bn{কিডনি} mean? (\textquotedblleft{}A kidney\textquotedblright{}.) Say it once more.
\end{tcolorbox}
\section[\texorpdfstring{\bn{ব্যায়াম} (bæyām)}{ব্যায়াম (bæyām)}]{\bn{ব্যায়াম} (bæyām) — exercise}

\label{lesson:BN-C155-byam}



\begin{quote}
Before the new one: say the Bengali for a lung, then the Bengali for a kidney.
\end{quote}

\subsection*{You'll want to know: \bn{ব্যায়াম}}
\textbf{\bn{ব্যায়াম}} — \emph{bæyām} — \textquotedblleft{}exercise\textquotedblright{}.

\begin{cousinweb}
\textbf{\bn{ব্যায়াম}} (\emph{bæyām}) is Sanskrit \emph{vyāyāma}, \textquotedblleft{}exertion\textquotedblright{}, taken over whole as a learned loan, a tatsama (CDIAL 12209).
\end{cousinweb}

\begin{grammarlens}[title={What you've built}]
One more word for this chapter.
\end{grammarlens}

\subsection*{Your turn}
\begin{itemize}
\raggedright
  \item \emph{Say it:} \emph{bæyām}
  \item \emph{Say it:} \emph{bæyām}, once more
  \item \emph{Say it:} say \emph{bæyām}
  \item \emph{Recall:} say the Bengali for a lung, then the Bengali for a kidney, then say \emph{bæyām} again
\end{itemize}

\begin{tcolorbox}[breakable,colback=teal!4,colframe=teal!35!black,title={Before you move on}]
What does \bn{ব্যায়াম} mean? (\textquotedblleft{}Exercise\textquotedblright{}.) Say it once more.
\end{tcolorbox}
\section[\texorpdfstring{\bn{স্বাস্থ্য} (svāsthyô)}{স্বাস্থ্য (svāsthyô)}]{\bn{স্বাস্থ্য} (svāsthyô) — health}

\label{lesson:BN-C155-svasthyo}



\begin{quote}
Before the new one: say the Bengali for a kidney, then the Bengali for exercise.
\end{quote}

\subsection*{You'll want to know: \bn{স্বাস্থ্য}}
\textbf{\bn{স্বাস্থ্য}} — \emph{svāsthyô} — \textquotedblleft{}health\textquotedblright{}.

\begin{cousinweb}
\textbf{\bn{স্বাস্থ্য}} (\emph{svāsthyô}) is Sanskrit \emph{svāsthya}, \textquotedblleft{}health, well-being\textquotedblright{}, taken over whole as a learned loan, a tatsama.
\end{cousinweb}

\begin{grammarlens}[title={What you've built}]
One more word for this chapter.
\end{grammarlens}

\subsection*{Your turn}
\begin{itemize}
\raggedright
  \item \emph{Say it:} \emph{svāsthyô}
  \item \emph{Say it:} \emph{svāsthyô}, once more
  \item \emph{Say it:} say \emph{svāsthyô}
  \item \emph{Recall:} say the Bengali for a kidney, then the Bengali for exercise, then say \emph{svāsthyô} again
\end{itemize}

\begin{tcolorbox}[breakable,colback=teal!4,colframe=teal!35!black,title={Before you move on}]
What does \bn{স্বাস্থ্য} mean? (\textquotedblleft{}Health\textquotedblright{}.) Say it once more.
\end{tcolorbox}
